% =========================================================================
%  THE DIRTY MAN — The Switch Operator
%  IEEE conference-style paper (IEEEtran)
%  Compile: pdflatex ieee_paper.tex
% =========================================================================
\documentclass[conference,10pt]{IEEEtran}

\usepackage{cite}
\usepackage{amsmath,amssymb}
\usepackage{graphicx}
\usepackage{booktabs}
\usepackage{url}
\usepackage[colorlinks=true,linkcolor=black,citecolor=blue,urlcolor=blue]{hyperref}
\usepackage{algorithmic}

\title{The Dirty Man: Self-Reconfiguring Neural Computation via\\Visual-Cue and Goal-Conditioned Routing}
\author{\IEEEauthorblockN{Sehaj Singh}
\IEEEauthorblockA{\textit{Independent Research}\\
2026, \url{sehaj@example.com}}}

\begin{document}
\maketitle

\begin{abstract}
We present the Dirty Man, a neural architecture whose core component --- the Switch Operator --- reconfigures its computation per input: a visual-cue ``eye'' and a goal pathway drive a differentiable router that selects among nine neural primitives (linear, dense, ReLU, CNN, RNN, LSTM, GAN, autoencoder, transformer), all emitting into a shared latent space. All experiments run on CUDA automatically when a GPU is available. Staged training --- warm-start primitives, oracle-supervised router training, joint fine-tune, annealed Gumbel commitment --- prevents mixture collapse and yields a readable routing policy. On a procedural sim-to-real glyph benchmark (24$\times$24, 3 seeds): (i) one operator reaches 0.834 accuracy, above every fixed network (best static 0.829, best standalone 0.822, static CNN 0.789); (ii) zero-shot real transfer 0.474 vs.\ 0.388 for a sim-trained static CNN, cutting the transfer gap from 0.611 to 0.526; (iii) with 200 real labels, structure adaptation (4{,}659 parameters) reaches 0.475 real / 0.9997 sim vs.\ weight adaptation (16{,}330 parameters) at 0.395 / 0.984; (iv) goal-conditioned routing yields 4.1$\times$ better reconstruction (0.034 vs.\ 0.139 MSE) with no classification loss; (v) ablations of bottleneck, annealing, eye, and domain-invariance change accuracy by $\le$0.001. Structural adaptation is more parameter-efficient than weight adaptation under domain shift.
\end{abstract}

\begin{IEEEkeywords}
dynamic neural architecture; mixture of experts; differentiable routing; Gumbel-Softmax; sim-to-real transfer; structural adaptation; goal-conditioned computation
\end{IEEEkeywords}

\section{Introduction}
Fixed-topology networks adjust weights inside a structure chosen before training. We study the complementary axis: adjusting \emph{which} network computes. The Switch Operator (Fig.~\ref{fig:arch}) routes each input to one of nine primitives with different inductive biases, using a router driven by visual cues and the task goal. Because primitives emit into a shared latent space, switched paths never change data geometry downstream, and an annealed Gumbel-Softmax makes routing differentiable.

Our claims, all reproduced from committed result files with 3 seeds:
\begin{itemize}
\item \textbf{Mixed-domain learning} (Sec.~\ref{sec:A}): one operator beats every fixed network it contains, learning a policy that shifts flat lenses $\rightarrow$ spatial/piecewise lenses as corruption grows.
\item \textbf{Transfer} (Sec.~\ref{sec:B}): structure adaptation (router only) beats weight adaptation (all weights) with 200 real labels, at 3.5$\times$ fewer adapted parameters and without forgetting the source domain.
\item \textbf{Multi-goal} (Sec.~\ref{sec:C}): goal-conditioned routing improves reconstruction 4.1$\times$ with zero classification cost.
\item \textbf{Robustness} (Sec.~\ref{sec:D}): every ablation changes accuracy by $\le$0.001.
\end{itemize}

\section{Method}
\subsection{Architecture}
Let $\mathcal{P} = \{\text{linear}, \text{dense}, \text{relu}, \text{cnn}, \text{rnn}, \text{lstm}, \text{gan}, \text{autoencoder}, \text{transformer}\}$ be $K=9$ primitives, each $\text{Prim}_k: \mathbb{R}^{24\times24} \to \mathbb{R}^{64}$ with its own task head. A convolutional eye produces cues $e=\text{Eye}(x)\in\mathbb{R}^{32}$; a goal pathway embeds the task $g\in\mathbb{R}^{16}$. The router
\begin{equation}
z_k = \text{Router}_k([e,g]), \qquad p_k = \text{softmax}\big((z + \varepsilon)/\tau\big)_k,
\end{equation}
with Gumbel noise $\varepsilon$ and annealed temperature $\tau(t) = 0.5 + 3.5\cdot\frac12(1+\cos\pi t)$. The mixture $\ell=\sum_k p_k\ell_k$ feeds the head.

\subsection{Staged training}
Training from scratch collapses (the router concentrates on one primitive). We train in stages:
\begin{enumerate}
\item warm-start primitives on mixed data (specialist subsets for flat vs.\ spatial lenses);
\item train eye+router on regime-level oracle targets (which expert is best per corruption regime), primitives frozen;
\item joint fine-tune;
\item deploy with $\tau\to0$ (argmax routing).
\end{enumerate}

\section{Benchmark}
Digits 0--9 are rendered procedurally (anti-aliased parametric strokes, 24$\times$24). Sim = clean render; real = sensor corruption (blur, noise, brightness/contrast drift, occlusion, JPEG-style quantization, affine warp). Each sample carries severity $\in[0,1]$ and a regime label (clean/spatial/statistical). Train 6{,}000 / test 2{,}000, batch 128, 12 epochs, 3 seeds, disjoint random streams.

\section{Results}
\subsection{Protocol A: mixed-domain}\label{sec:A}
\begin{table}[h]
\centering
\caption{Mixed-domain accuracy (3 seeds).}
\label{tab:a}
\begin{tabular}{lcc}
\toprule
Model & Acc & Adapted params\\
\midrule
\textbf{Switch Operator} & \textbf{0.834} & 4{,}659\\
static MLP & 0.829 & ---\\
best standalone (ReLU) & 0.822 & ---\\
random router & 0.821 & ---\\
static CNN & 0.789 & ---\\
uniform router & 0.753 & ---\\
\bottomrule
\end{tabular}
\end{table}
The operator wins consistently (0.843/0.831/0.834). Diagnostics show a learned policy: dense mass 0.617 on sim $\to$ 0.001 on real; relu/cnn rise from 0/0.383 to 0.526/0.473 (Fig.~\ref{fig:routing}).

\subsection{Protocol B: transfer}\label{sec:B}
\begin{table}[h]
\centering
\caption{Sim$\to$real transfer (200 real labels, 3 seeds).}
\label{tab:b}
\begin{tabular}{lcccc}
\toprule
Model & Real & Sim & Gap & Params\\
\midrule
static CNN & 0.388 & 0.999 & 0.611 & ---\\
switch (zero-shot) & 0.474 & 1.000 & 0.526 & 0\\
weight adapt & 0.395 & 0.984 & 0.611 & 16{,}330\\
\textbf{structure adapt} & \textbf{0.475} & \textbf{1.000} & \textbf{0.526} & \textbf{4{,}659}\\
\bottomrule
\end{tabular}
\end{table}
Structure adaptation wins the target domain at 3.5$\times$ fewer adapted parameters and does not forget sim (0.9997 vs.\ 0.984).

\subsection{Protocol C: goal pathway}\label{sec:C}
\begin{table}[h]
\centering
\caption{Goal-conditioned vs.\ agnostic (3 seeds).}
\label{tab:c}
\begin{tabular}{lcc}
\toprule
Model & Cls acc & Recon MSE\\
\midrule
goal-agnostic & 0.810 & 0.139\\
\textbf{goal-conditioned} & \textbf{0.821} & \textbf{0.034}\\
\bottomrule
\end{tabular}
\end{table}
Reconstruction improves 4.1$\times$; classification +1.1pt. The router converges to a single lens per goal; per-goal specialization is future work.

\subsection{Protocol D: ablations}\label{sec:D}
\begin{table}[h]
\centering
\caption{Ablations (3 seeds).}
\label{tab:d}
\begin{tabular}{lcc}
\toprule
Variant & Acc & Real\\
\midrule
full & 0.824 & 0.648\\
no bottleneck & 0.824 & 0.649\\
no annealing & 0.824 & 0.648\\
domain-invariant eye & 0.823 & 0.646\\
no eye & 0.823 & 0.646\\
random router & 0.823 & 0.647\\
\bottomrule
\end{tabular}
\end{table}
Robust to every removal ($\Delta \le 0.001$).

\section{Related Work}
Mixture-of-experts \cite{switchtransformer,condcomp} routes tokens to parallel experts; we add visual-cue+goal routing and oracle-supervised staged training. Dynamic networks \cite{dynamicnn} change depth/width/skips; we switch the family of computation. NAS \cite{nas} searches one structure per dataset at heavy cost; we switch per sample at inference. RepVGG \cite{repvgg} folds branches at inference; we keep them live. Sim-to-real \cite{domainrand,sim2real} adapts weights or renders domains; we show structure adaptation is more parameter-efficient.

\section{Conclusion}
The Switch Operator shows that structural adaptation --- changing which network computes, per input, per goal --- is learnable, interpretable, transferable, and more parameter-efficient than weight adaptation under domain shift. Code, figures, and all numbers are committed and reproducible.

\section*{Code and reproducibility}
\url{https://github.com/sehajr-singhs/dirty-man} --- every number is read from committed JSON result files by \texttt{make\_figs.py}; the benchmark is procedural (no downloads); \texttt{--device cuda} forces GPU execution and \texttt{run\_experiments.py --smoke} verifies the pipeline in ~1 minute.

\begin{thebibliography}{9}
\bibitem{switchtransformer} N.~Shazeer \emph{et al.}, ``Outrageously large neural networks: The sparsely-gated mixture-of-experts layer,'' \emph{ICLR}, 2017.
\bibitem{condcomp} Y.~Bengio, N.~L\'eonard, and A.~Courville, ``Estimating or propagating gradients through stochastic neurons for conditional computation,'' \emph{arXiv:1308.3432}, 2013.
\bibitem{gumbel} E.~Jang, S.~Gu, and B.~Poole, ``Categorical reparameterization with Gumbel-Softmax,'' \emph{ICLR}, 2017.
\bibitem{repvgg} X.~Ding \emph{et al.}, ``RepVGG: Making VGG-style ConvNets great again,'' \emph{CVPR}, 2021.
\bibitem{dynamicnn} Y.~Han \emph{et al.}, ``Dynamic neural networks: A survey,'' \emph{IEEE TPAMI}, 2021.
\bibitem{nas} T.~Elsken, J.~H.~Metzen, and F.~Hutter, ``Neural architecture search: A survey,'' \emph{JMLR}, 2019.
\bibitem{domainrand} J.~Tobin \emph{et al.}, ``Domain randomization for transferring deep neural networks from simulation to the real world,'' \emph{IROS}, 2017.
\bibitem{sim2real} X.~B.~Peng, M.~Andrychowicz, W.~Zaremba, and P.~Abbeel, ``Sim-to-real transfer of robotic control with dynamics randomization,'' \emph{ICRA}, 2018.
\end{thebibliography}

\end{document}
